\documentclass[conference,a4paper]{IEEEtran}

\usepackage[T1]{fontenc}
\usepackage[utf8]{inputenc}
\usepackage{amsmath}
\usepackage{microtype}

\begin{document}

\title{Integrated Pick Scheduling and Conveyor Speed Control for a
Delta-Robot Sorting Cell}

\author{
\IEEEauthorblockN{Tran Manh Son, Dao Minh Thong, Tran Dang Kim Toan}
}

\maketitle

\begin{abstract}
This study develops an online algorithm that jointly schedules picks and
controls conveyor speed for a delta-robot sorting system. Each part on the
conveyor stays within the robot's reachable area only for a limited time; if
it is not picked before leaving, it is lost. The goal is to pick as many
parts as possible before time runs out -- an online scheduling problem with
time-varying release times and deadlines, $1 \mid r_j(t), d_j(t),
\text{online} \mid \sum U_j$. These deadlines are not fixed: they depend on
the conveyor speed, which the controller can adjust. Slowing the belt gives
more time for the part being chased, but delays every part behind it. The
algorithm treats scheduling and conveyor speed as one problem, using a
two-loop structure. Because the camera sits upstream, the inner loop sees
each part before it arrives, decides which part to pick next, and predicts
how many parts on the belt will likely be missed. The outer loop uses that
prediction to adjust the belt speed. The time to switch between two picks
depends only on the type of the part just picked, so it is folded into the
pick time. The algorithm is validated on a self-built delta-robot cell --
vacuum gripper, rotation axis, encoder-equipped conveyor, a YOLO-OBB vision
model, two PLCs, and a control computer -- running at 30-60 picks per minute.
\end{abstract}

\begin{IEEEkeywords}
Delta robot; Online scheduling; Conveyor speed control; Sorting system;
Machine vision.
\end{IEEEkeywords}

\end{document}
